\section{PCB layout (pending)}\label{sec:layout}

This chapter will be filled in after layout. The inputs are ready:
\begin{itemize}
  \item the board envelopes and height zones (\cref{fig:mainboard-env,fig:uiboard-env});
  \item the layout checklists in both specs;
  \item the release reviews' placement conditions (tall capacitors, XT30 polarity, AGND/PGND tie);
  \item the project's layout principles, \repo{docs/layout_principles.md}.
\end{itemize}
The principles say that layout is where LLMs are weakest, so the board will be judged by physics, not just by connectivity and DRC:
\begin{itemize}
  \item The model writes net classes and rules from the datasheets' layout sections, with page citations.
  \item An autorouter handles non-critical nets.
  \item Power stages, switch nodes, the USB pair and high-current paths are routed deliberately and reviewed.
  \item The result is scored with loop-inductance, PDN and IR-drop checks. Post-layout parasitics go back into the S7 PDN model.
\end{itemize}

\begin{figure}[htbp]
\centering
\begin{subfigure}[t]{0.49\linewidth}\centering\pendingbox[0.95\linewidth]{4.5cm}\caption{Main board: placement and top copper.}\end{subfigure}\hfill
\begin{subfigure}[t]{0.49\linewidth}\centering\pendingbox[0.95\linewidth]{4.5cm}\caption{Main board: 3D view in the enclosure envelope.}\end{subfigure}
\caption{PCB layout: placeholders until the layout exists.}
\label{fig:layout-pending}
\end{figure}
